\documentclass[runningheads]{llncs}

% ===== Packages =====
\usepackage[T1]{fontenc}
\usepackage[utf8]{inputenc}
\usepackage[spanish]{babel}
\usepackage{amsmath,amssymb,amsfonts}
\usepackage{graphicx}
\usepackage{booktabs}
\usepackage{tikz}
\usetikzlibrary{arrows.meta,positioning}
\usetikzlibrary{babel}
\usepackage{algorithm}
\usepackage{algpseudocode}
\usepackage[hidelinks]{hyperref}

% ===== Math shortcuts =====
\newcommand{\kk}{\Bbbk}
\newcommand{\kQ}{\kk Q}
\newcommand{\NN}{\mathbb{N}}
\newcommand{\RR}{\mathbb{R}}
\newcommand{\Adj}{A}
\newcommand{\walk}{\rightsquigarrow}
\DeclareMathOperator{\softmax}{softmax}

\begin{document}

\title{Álgebras de Caminos de Quivers como una\\ Atención Dispersa Multi-Salto para Mitigar\\ el Over-Squashing en GNNs}
\titlerunning{Walk Attention}

\author{Carlos Isaac Zainea Maya\inst{1}\orcidID{0000-0002-6459-2397} \and
Agustín Moreno Cañadas\inst{1}\orcidID{0000-0001-6812-5131}}
\authorrunning{C. I. Zainea Maya and A. Moreno Cañadas}
\institute{Universidad Nacional de Colombia, Departamento de Matemáticas,\\
Bogotá, Colombia\\
\email{\{cizaineam,amorenoca\}@unal.edu.co}}

\maketitle

\begin{abstract}
Las redes neuronales de grafos de paso de mensajes sufren de
\emph{over-squashing}: la señal que viaja por muchos caminos largos hacia un
nodo se comprime en un vector de ancho fijo y se pierde. Introducimos
\textbf{Walk Attention}: el \emph{álgebra de caminos} $\kQ$ del quiver
subyacente provee, mediante sus componentes graduadas, la estructura exacta de
alcanzabilidad multi-salto, que usamos como \emph{soporte} disperso de una
atención aprendida---el álgebra dice qué vértices pueden interactuar a alcance
$k$, la red aprende cuánto. En grafos de cuello de botella con over-squashing
ajustable \emph{demostramos} la separación---el paso de mensajes con
agregación lineal queda exactamente en el azar para toda profundidad y todo
ancho, mientras que una sola capa de Walk Attention en el alcance adecuado
resuelve la tarea de recuperación exactamente---y la confirmamos
empíricamente ($1.000\pm0.000$ sobre cinco semillas, donde la atención de un
salto se degrada hacia el azar); Walk Attention se mantiene exacta en
\textsc{RingTransfer} y, bajo un presupuesto pequeño compartido, empata con la
atención de un salto en Peptides-func: su ventaja es decisiva bajo
over-squashing severo y modesta cuando es leve. Un resultado negativo afina el
diseño: colapsar caminos equivalentes vía el cociente $\kQ/I$ descarta
multiplicidad que la tarea necesita; el papel del álgebra es \emph{definir el
soporte}, no comprimir.

\keywords{Redes neuronales de grafos \and Over-squashing \and Álgebra de
caminos \and Quiver \and Teoría de representaciones \and Atención.}
\end{abstract}

% =====================================================================
\section{Introducción}
% =====================================================================

Las redes neuronales de grafos construidas sobre el paradigma de paso de
mensajes \cite{gilmer2017neural} actualizan cada nodo agregando mensajes de sus
vecinos inmediatos e iterando esta regla local a través de las capas, de modo
que tras $k$ capas la información de nodos a $k$ saltos alcanza un nodo
objetivo. Aunque simple y efectivo, el esquema tiene un modo de fallo bien
documentado: el \emph{over-squashing} \cite{alon2021bottleneck}. Cuando el
número de caminos que llevan información hacia un nodo crece---típicamente de
forma exponencial con la distancia---los mensajes a lo largo de todos ellos
deben comprimirse en un único vector de ancho fijo; el nodo receptor no puede
recuperarlos individualmente y la señal de largo alcance se pierde. El
over-squashing es el obstáculo principal para tareas cuyas etiquetas dependen
de información distante o distribuida globalmente. Los remedios existentes
actúan sobre el grafo o sobre la agregación: el \emph{rewiring} ensancha los
cuellos de botella \cite{topping2022understanding}, los análisis de curvatura y
sensibilidad explican \emph{dónde} ocurre \cite{digiovanni2023oversquashing},
y una línea complementaria reemplaza la agregación local por atención global de
Transformer \cite{vaswani2017attention}, a costo cuadrático y ciego a la
estructura del grafo.

En este artículo tomamos una ruta diferente, fundamentada en la teoría de
representaciones de quivers. Un \emph{quiver} es un grafo dirigido,
posiblemente con flechas múltiples \cite{schiffler2014quiver,assem2006elements},
y su \emph{álgebra de caminos} $\kQ$ tiene como base los recorridos dirigidos
del quiver, con la concatenación como multiplicación. La pieza graduada
$e_j\cdot\kQ_k\cdot e_i$ de esta álgebra tiene como base los recorridos
dirigidos de longitud~$k$ del vértice~$i$ al vértice~$j$; su dimensión es
$(\Adj^k)_{ij}$, la entrada correspondiente de la $k$-ésima potencia de la
matriz de adyacencia. El álgebra de caminos provee, por tanto, una
\emph{descripción exacta y algebraica de la alcanzabilidad y la multiplicidad
multi-salto}.

Nuestra observación central es que esta estructura es precisamente lo que se
necesita para definir una atención multi-salto dispersa con fundamento. Dejamos
que cada nodo atienda directamente a todo vértice alcanzable desde él por un
recorrido de longitud~$k$: el álgebra provee el \emph{soporte} de la atención
(qué pares pueden interactuar a alcance~$k$), y la red aprende los \emph{pesos}
(cuánto). Como el soporte es disperso---típicamente una pequeña fracción de
todos los pares---esto es mucho más barato que la auto-atención completa, sin
dejar de actuar a distancia, evitando el cuello de botella iterativo. Llamamos
a la arquitectura resultante \textbf{Walk Attention}. Como una red de grafos sobre
$Q$ es una representación del quiver
($\mathrm{Rep}(Q)\cong\kQ\text{-}\mathrm{Mod}$), Walk Attention es la lectura
\emph{aprendida} de un sustrato algebraico cuya lectura \emph{dinámica} es el
marco de autómatas de impacto sobre quivers \cite{zainea2026aiq}
(Sección~\ref{sec:ground}).

\paragraph{Contribuciones.}
\begin{enumerate}
  \item Una construcción autocontenida (Sección~\ref{sec:algebra}) de quivers,
  álgebras de caminos, el sistema de vecindades por grado de impacto y los
  \emph{operadores de recorrido} que codifican la alcanzabilidad multi-salto,
  con ejemplos desarrollados.
  \item \textbf{Walk Attention} (Sección~\ref{sec:model}): una atención
  multi-salto aprendida cuyo soporte es la alcanzabilidad del álgebra de
  caminos, relacionada con el paso de mensajes y con los Transformers.
  \item Una separación \emph{demostrada} (Sección~\ref{sec:theory}): con la
  profundidad estándar de una capa por salto, toda GNN cuya primera agregación
  es lineal en los atributos de las fuentes (GCN, GIN) tiene precisión
  exactamente igual al azar en la familia de cuellos de botella, para toda
  profundidad, ancho y parámetros, mientras que una sola capa de Walk Attention
  en el alcance adecuado resuelve la tarea exactamente---confirmado
  empíricamente (Sección~\ref{sec:exp}).
  \item Un resultado \emph{negativo} que afina el diseño: colapsar caminos
  equivalentes vía el cociente $\kQ/I$ descarta multiplicidad necesaria y
  falla; el papel del cociente es determinar el soporte, no comprimir la señal.
\end{enumerate}

% =====================================================================
\section{Trabajo Relacionado}\label{sec:related}
% =====================================================================

\paragraph{Over-squashing, atención y agregación multi-salto.}
Alon y Yahav \cite{alon2021bottleneck} identificaron el over-squashing con la
prueba \textsc{NeighborsMatch}; Topping et al.\ \cite{topping2022understanding}
lo conectaron con la curvatura negativa y propusieron el rewiring SDRF; Di
Giovanni et al.\ \cite{digiovanni2023oversquashing} dieron un análisis de
sensibilidad.
GAT \cite{velickovic2017graph} aprende pesos de atención de un salto; los graph
Transformers aplican atención de todos los pares \cite{vaswani2017attention} a
costo $O(n^2)$.
Agregar más allá de un salto está, en sí, bien explorado: MixHop
\cite{abu2019mixhop} y SIGN \cite{frasca2020sign} mezclan potencias fijas de la
adyacencia, la convolución por difusión \cite{gasteiger2019diffusion} recablea
con un núcleo de difusión, MAGNA \cite{wang2021magna} difunde los puntajes de
atención, el paso de mensajes por caminos más cortos agrega sobre capas de
distancia \cite{abboud2022shortest}, NAGphormer \cite{chen2023nagphormer}
tokeniza agregados por salto, Graphormer \cite{ying2021graphormer} sesga la
atención densa por la distancia más corta, y DRew \cite{gutteridge2023drew}
recablea dinámicamente con retardo. Walk Attention difiere en el objeto que
define el soporte: el \emph{álgebra de caminos graduada}---recorridos con su
multiplicidad, no capas de distancia ni pesos de difusión---lo que hace el
soporte exacto (Proposición~\ref{prop:count}), expone el cociente $\kQ/I$ como
grado de libertad de diseño e inserta la arquitectura en la teoría de
representaciones de quivers.

\paragraph{Estructura algebraica de caminos.}
El álgebra de caminos es un objeto estándar de la teoría de representaciones
\cite{schiffler2014quiver,assem2006elements}; las álgebras de configuración de
Brauer y construcciones afines se han aplicado recientemente a problemas
combinatorios por Moreno Cañadas y colaboradores
\cite{moreno2021brauer,moreno2024integer}. Lo más cercano al presente trabajo
es el sistema de vecindades por grado de impacto y el marco de autómatas de
impacto \cite{zainea2026aiq,ibanez2025casimulations}, que tomamos como el
hermano dinámico de Walk Attention sobre el mismo sustrato
(Sección~\ref{sec:ground}). La atención sobre soportes multi-salto no es nueva
en sí misma; lo nuevo, hasta donde sabemos, es tomar el soporte del álgebra de
caminos graduada, con sus multiplicidades, su cociente y su lectura dinámica.

% =====================================================================
\section{El Álgebra de Caminos de un Quiver}\label{sec:algebra}
% =====================================================================

Esta sección es autocontenida: construimos los quivers, los caminos, el álgebra
de caminos con su graduación, el sistema de vecindades por grado de impacto y
los operadores de recorrido en que se apoya el modelo de la
Sección~\ref{sec:model}, sin suponer exposición previa. Referencias estándar:
\cite{schiffler2014quiver,assem2006elements}.

% ---------------------------------------------------------------------
\subsection{Quivers y caminos}
% ---------------------------------------------------------------------

\begin{definition}
Un \emph{quiver} es una cuádrupla $Q=(Q_0,Q_1,s,t)$ donde $Q_0$ es un conjunto
finito de \emph{vértices}, $Q_1$ un conjunto finito de \emph{flechas} y
$s,t\colon Q_1\to Q_0$ asignan a cada flecha $\alpha$ su fuente $s(\alpha)$ y
su destino $t(\alpha)$; escribimos $\alpha\colon i\to j$. Se permiten flechas
múltiples entre el mismo par ordenado. Identificamos $Q_0$ con $\{1,\dots,n\}$.
\end{definition}

Un quiver es así un multigrafo dirigido, el dato de una estructura relacional
dirigida: entidades, relaciones dirigidas y flechas paralelas que registran
multiplicidad.

\begin{example}\label{ex:firstquiver}
Sea $Q_0=\{1,2,3\}$ con flechas $\alpha,\alpha'\colon 1\to 2$ (dos flechas
paralelas), $\beta\colon 2\to 3$ y un lazo $\gamma\colon 3\to 3$. Su matriz de
adyacencia, contando flechas, es
\[
  \Adj=\begin{pmatrix}0&2&0\\0&0&1\\0&0&1\end{pmatrix}.
\]
En una red de citaciones, por ejemplo, los vértices son documentos, las flechas
citaciones, el par paralelo $\alpha,\alpha'$ una citación doble y el lazo
$\gamma$ una autocitación.
\end{example}

\begin{definition}
Un \emph{camino de longitud} $k\ge1$ en $Q$ es una secuencia de flechas
$p=\alpha_k\cdots\alpha_2\alpha_1$ con $t(\alpha_r)=s(\alpha_{r+1})$, leída de
derecha a izquierda como la composición de funciones, con fuente
$s(p)=s(\alpha_1)$ y destino $t(p)=t(\alpha_k)$. Cada vértice $i$ porta además
un \emph{camino trivial} $e_i$ de longitud $0$.
\end{definition}

Un camino es un \emph{recorrido} dirigido: vértices y flechas pueden repetirse.
En el Ejemplo~\ref{ex:firstquiver}, $\beta\alpha$ y $\beta\alpha'$ son dos
caminos distintos de longitud $2$ de $1$ a $3$, y el lazo genera recorridos
arbitrariamente largos $\gamma^m\beta\alpha$. El número de recorridos distintos
de una longitud dada entre dos vértices es precisamente la cantidad que produce
el over-squashing, y el álgebra de caminos lo hace algebraicamente explícito.

% ---------------------------------------------------------------------
\subsection{El álgebra de caminos y su graduación}
% ---------------------------------------------------------------------

Sea $\kk$ un cuerpo (en la práctica $\kk=\RR$).

\begin{definition}
El \emph{álgebra de caminos} $\kQ$ es el $\kk$-espacio vectorial con base el
conjunto de todos los caminos de $Q$, dotado de la multiplicación bilineal dada
en la base por la concatenación:
\[
  q\cdot p \;=\;
  \begin{cases}
    q\,p & \text{si } s(q)=t(p),\\[2pt]
    0 & \text{en otro caso.}
  \end{cases}
\]
Los caminos triviales satisfacen $e_{t(p)}\,p=p=p\,e_{s(p)}$ y
$e_ie_j=\delta_{ij}e_i$, de modo que $\{e_i\}$ es un sistema completo de
idempotentes ortogonales y $1=\sum_i e_i$ es la unidad.
\end{definition}

La multiplicación respeta la longitud: concatenar caminos de longitudes $k$ y
$\ell$ da un camino de longitud $k+\ell$ (o cero). Por tanto $\kQ$ es una
álgebra \emph{graduada},
\[
  \kQ \;=\; \bigoplus_{k\ge 0}\kQ_k,
  \qquad \kQ_k\cdot \kQ_\ell \subseteq \kQ_{k+\ell},
\]
con $\kQ_k$ generada por los caminos de longitud exactamente $k$; en particular
$\kQ_0=\bigoplus_i\kk e_i$ y $\kQ_1$ está generada por las flechas. Los
idempotentes permiten aislar los caminos entre un par fijo de vértices: el
subespacio $e_j\cdot\kQ_k\cdot e_i$ está generado exactamente por los caminos
de longitud $k$ de $i$ a $j$, y su dimensión es por tanto el número de tales
recorridos. La siguiente proposición, en la que se apoya todo lo que el modelo
calcula, hace precisa la correspondencia con el álgebra lineal.

\begin{proposition}\label{prop:count}
Sea $\Adj\in\NN^{n\times n}$ la matriz de adyacencia de $Q$,
$\Adj_{ij}=\#\{\alpha\in Q_1:\alpha\colon i\to j\}$. Entonces, para todo
$k\ge0$ y todos $i,j\in Q_0$,
\[
  \dim_{\kk}\bigl(e_j\cdot \kQ_k\cdot e_i\bigr) \;=\; (\Adj^{\,k})_{ij}.
\]
\end{proposition}

\begin{proof}
Inducción sobre $k$; para $k\in\{0,1\}$ es la definición. Todo camino de
longitud $k$ de $i$ a $j$ se factoriza de forma única como $\alpha\cdot p$, con
$p$ de longitud $k-1$ de $i$ a algún vértice intermedio $r$ y
$\alpha\colon r\to j$ una flecha; el número de tales factorizaciones es
$\sum_r(\Adj^{\,k-1})_{ir}\Adj_{rj}=(\Adj^{\,k})_{ij}$. \qed
\end{proof}

Llamamos a $n_k(i,j):=(\Adj^{\,k})_{ij}=\dim_\kk(e_j\kQ_k e_i)$ la
\emph{multiplicidad de recorrido} de $i$ a $j$ a alcance $k$. Es exactamente el
número de mensajes de longitud $k$ que una GNN de paso de mensajes fuerza del
nodo $i$ hacia el nodo $j$ tras $k$ capas---y por tanto la fuente del
over-squashing.

\begin{example}\label{ex:diamond}
Sea $Q$ con vértices $\{1,2,3,4\}$ y flechas $\alpha_1\colon1\to2$,
$\alpha_2\colon1\to3$, $\beta_1\colon2\to4$, $\beta_2\colon3\to4$
(Figura~\ref{fig:diamond}). Entonces
\[
  \Adj=\begin{pmatrix}0&1&1&0\\0&0&0&1\\0&0&0&1\\0&0&0&0\end{pmatrix},
  \qquad
  \Adj^2=\begin{pmatrix}0&0&0&2\\0&0&0&0\\0&0&0&0\\0&0&0&0\end{pmatrix},
\]
y $\dim(e_4\kQ_2e_1)=(\Adj^2)_{14}=2$, de acuerdo con la
Proposición~\ref{prop:count}: los dos recorridos \emph{paralelos} de longitud
$2$, $\beta_1\alpha_1$ y $\beta_2\alpha_2$, con la misma fuente, el mismo
destino y la misma longitud. Si deben tratarse como un canal o como dos es una
decisión de modelado a la que volvemos en la Sección~\ref{sec:exp}.
\end{example}

\begin{figure}[t]
\centering
\begin{tikzpicture}[->,>={Stealth[round]},node distance=8mm and 12mm,
  every node/.style={circle,draw,minimum size=6mm,inner sep=0pt,font=\scriptsize}]
  \node (1) {1};
  \node (2) [above right=of 1] {2};
  \node (3) [below right=of 1] {3};
  \node (4) [below right=of 2] {4};
  \draw (1) -- node[draw=none,fill=none,above left,font=\tiny]{$\alpha_1$} (2);
  \draw (1) -- node[draw=none,fill=none,below left,font=\tiny]{$\alpha_2$} (3);
  \draw (2) -- node[draw=none,fill=none,above right,font=\tiny]{$\beta_1$} (4);
  \draw (3) -- node[draw=none,fill=none,below right,font=\tiny]{$\beta_2$} (4);
\end{tikzpicture}
\caption{El quiver diamante del Ejemplo~\ref{ex:diamond}: dos recorridos
paralelos de longitud $2$, $1\walk4$; $\dim(e_4\kQ_2 e_1)=2$.}
\label{fig:diamond}
\end{figure}

% ---------------------------------------------------------------------
\subsection{Grado de impacto y el sistema fundamental de vecindades}
\label{sec:sfv}
% ---------------------------------------------------------------------

La graduación por longitud induce una estratificación canónica de los vértices
alrededor de cada nodo, que organiza la estructura multi-salto que el modelo
explotará. El \emph{grado de impacto} $g(i,j)$ es la longitud del camino
dirigido más corto de $i$ a $j$ (con $g(i,i)=0$ y $g(i,j)=\infty$ si no
existe); equivalentemente, $g(i,j)=\min\{k:(\Adj^{\,k})_{ij}>0\}$, el primer
alcance en que la multiplicidad se hace positiva.

\begin{definition}\label{def:sfv}
Para un nodo $c\in Q_0$, su \emph{vecindad de entrada} de radio $k$ es
$A_k(c)=\{i\in Q_0: g(i,c)\le k\}$, el conjunto de nodos que pueden influir en
$c$ por un camino dirigido de longitud a lo sumo $k$. La familia anidada
$\{A_k(c)\}_{k\ge0}$, con $A_0(c)=\{c\}\subseteq A_1(c)\subseteq\cdots$, es el
\emph{sistema fundamental de vecindades} (SFV) de $c$; la capa
$\mathcal{N}_k(c)=A_k(c)\setminus A_{k-1}(c)$ reúne los nodos a grado de
impacto exactamente $k$.
\end{definition}

El SFV es el objeto combinatorio sobre el que actúan los operadores de
recorrido de más abajo, y su relación precisa con la alcanzabilidad es la
siguiente.

\begin{proposition}\label{prop:fns}
Para $k\ge1$ sea $\mathcal{S}_k=\{(c,i):(\Adj^{\,k})_{ic}>0\}$ el conjunto de
pares unidos por un recorrido de longitud $k$, $i\walk c$. Entonces
\[
  \{(c,i): i\in \mathcal{N}_k(c)\} \;\subseteq\; \mathcal{S}_k ,
\]
y la inclusión es estricta en general: los recorridos de longitud $k$ pueden
alcanzar también nodos a grado de impacto $<k$. Si $Q$ es
\emph{graduado}---existe $\rho\colon Q_0\to\mathbb{Z}$ con
$\rho(t(\alpha))=\rho(s(\alpha))+1$ para toda flecha $\alpha$---vale la
igualdad para todo $k$.
\end{proposition}

\begin{proof}
Si $g(i,c)=k$, un camino más corto es un recorrido de longitud $k$, así que
$(\Adj^{\,k})_{ic}>0$. Estrictez: en $1\to2\to3$ con una flecha adicional
$1\to3$, $(\Adj^2)_{13}>0$ y sin embargo $g(1,3)=1$. Si $Q$ es graduado, todo
recorrido $i\walk c$ tiene longitud $\rho(c)-\rho(i)$, así que un recorrido de
longitud $k$ fuerza $g(i,c)=k$. \qed
\end{proof}

Los grafos del benchmark de la Sección~\ref{sec:exp} son graduados (el índice
de capa es un tal $\rho$), así que allí los soportes coinciden con las capas
del SFV; en grafos con ciclos, como un anillo, $\mathcal{S}_k$ contiene
estrictamente la relación de capas. La multiplicidad $n_k(i,c)$ registra
\emph{cuántos} caminos de longitud $k$ pueblan cada capa---el refinamiento
algebraico del grado de impacto escalar.

% ---------------------------------------------------------------------
\subsection{Operadores de recorrido}
% ---------------------------------------------------------------------

La Proposición~\ref{prop:count} permite convertir el álgebra graduada en una
familia de operadores lineales sobre los atributos, uno por alcance; son el
puente hacia el modelo neuronal.

\begin{definition}\label{def:walkop}
Para cada alcance $k\ge1$ defínase $W^{(k)}\in\RR^{n\times n}$ por
\[
  W^{(k)}_{j i}\;=\;\frac{(\Adj^{\,k})_{ij}}{\,Z_k\,},
  \qquad Z_k=\max_{j}\sum_i (\Adj^{\,k})_{ij},
\]
una sola constante de normalización por alcance. El \emph{soporte de
alcanzabilidad por recorridos} es
$\mathcal{S}_k=\{(j,i):(\Adj^{\,k})_{ij}>0\}$.
\end{definition}

Actuando sobre una matriz de atributos $H\in\RR^{n\times d}$, el operador
$W^{(k)}H$ envía a cada nodo $j$ una suma ponderada de los atributos de cada
nodo $i$ que lo alcanza por un recorrido de longitud $k$, con peso proporcional
a la multiplicidad $n_k(i,j)$. Dos hechos merecen énfasis.
\emph{(i) Normalización global, no por filas}: normalizar por filas mapearía
las multiplicidades a una distribución uniforme sobre las fuentes alcanzables,
borrando exactamente la información que distingue a $W^{(k)}$ de una simple
máscara de alcanzabilidad; la constante global preserva las multiplicidades
relativas acotando las magnitudes. \emph{(ii) Sparsity}: $\mathcal{S}_k$ es el
patrón de no-ceros de $\Adj^{\,k}$, así que en los grafos de interés $W^{(k)}$
se almacena y aplica en $O(|\mathcal{S}_k|)$ y no en $O(n^2)$. El operador de
pesos fijos $W^{(k)}$ ya mitiga el over-squashing, porque deja que el nodo $j$
lea las fuentes a alcance $k$ \emph{directamente}, en un paso, en vez de forzar
la señal por $k$ rondas locales. Su limitación---que motiva Walk Attention---es
que sus pesos están fijados por la multiplicidad y no pueden \emph{seleccionar}
qué fuente importa.

% ---------------------------------------------------------------------
\subsection{El álgebra cociente}
% ---------------------------------------------------------------------

La teoría de representaciones pasa rutinariamente de $\kQ$ a un cociente
$\kQ/I$ por un \emph{ideal admisible} $I$, identificando caminos declarados
iguales. En el diamante del Ejemplo~\ref{ex:diamond}, la relación
$\beta_1\alpha_1\equiv\beta_2\alpha_2$ genera un ideal con
$\dim\bigl(e_4(\kQ/I)_2e_1\bigr)=1$: los dos recorridos paralelos colapsan a
una sola clase, y en general
$\dim\bigl(e_j(\kQ/I)_ke_i\bigr)\le n_k(i,j)$. Es tentador usarlo como remedio
del over-squashing: reemplazar las multiplicidades por las dimensiones menores,
deduplicando caminos antes de agregar. Lo probamos y resultó
\emph{contraproducente}
(Sección~\ref{sec:exp}): colapsar recorridos paralelos destruye información de
multiplicidad que las tareas pueden requerir. El papel constructivo del
cociente no es comprimir la señal sino \emph{remodelar el soporte}---por
ejemplo enrutando clases de caminos distintas a cabezas de atención distintas,
dirección que dejamos a trabajo futuro. En el modelo que sigue el álgebra entra
solo por el soporte $\mathcal{S}_k$ y los operadores sin cociente $W^{(k)}$.

% ---------------------------------------------------------------------
\subsection{Un suelo algebraico común}\label{sec:ground}
% ---------------------------------------------------------------------

Los objetos anteriores---el álgebra graduada $\kQ$, el grado de impacto con su
SFV y el cociente $\kQ/I$---forman un sustrato algebraico que se especializa en
dos direcciones, según cómo se lea un mapa de actualización $\varphi$ sobre las
capas $\mathcal{N}_k(c)$. En la lectura \emph{dinámica}, $\varphi$ es una regla
de evolución fija ponderada por el grado de impacto; iterarla en el tiempo da
los \emph{autómatas de impacto sobre quivers} (AIQ), con sus reglas SI/SIS/SIR
y su estructura categórica, desarrollados aparte
\cite{zainea2026aiq,ibanez2025casimulations}. En la lectura
\emph{representacional}, tomada aquí, $\varphi$ se \emph{aprende}: un espacio
vectorial por vértice y un mapa lineal por flecha es una representación de $Q$
(un $\kQ$-módulo, vía $\mathrm{Rep}(Q)\cong\kQ\text{-}\mathrm{Mod}$), y una GNN
sobre $Q$ es tal representación con mapas aprendidos. Walk Attention es su
instancia graduada.

% =====================================================================
\section{Walk Attention}\label{sec:model}
% =====================================================================

Definimos ahora la arquitectura, y empezamos por la observación que la motiva:
\emph{el operador de recorrido ya es una atención}. Una capa de atención, en su
forma general, fija una relación de \emph{soporte} $S$ de pares (objetivo,
fuente) y calcula, para cada objetivo $j$, una combinación ponderada y
normalizada de atributos transformados de las fuentes,
\begin{equation}\label{eq:attn-generic}
  \mathrm{out}_j \;=\; \sum_{i:\,(j,i)\in S}\alpha_{ji}\,Vx_i,
  \qquad \alpha_{ji}\ge0,\quad \textstyle\sum_{i}\alpha_{ji}\;\text{acotada},
\end{equation}
y las arquitecturas difieren solo en dos elecciones: \emph{qué pares} contiene
$S$ y \emph{de dónde salen los coeficientes} $\alpha$. La auto-atención de los
Transformers \cite{vaswani2017attention} toma $S=$ todos los pares y aprende
$\alpha$ de productos consulta--clave; GAT \cite{velickovic2017graph} toma $S=$
las aristas y aprende $\alpha$. Compárese ahora con el operador de la
Definición~\ref{def:walkop}: actuando sobre atributos calcula
$(W^{(k)}H)_j=\sum_{i:(j,i)\in\mathcal{S}_k}\tfrac{n_k(i,j)}{Z_k}\,H_i$---
\emph{exactamente la forma}~\eqref{eq:attn-generic}, con soporte la
alcanzabilidad $\mathcal{S}_k$ y coeficientes fijados en las multiplicidades
normalizadas. Agregar sobre recorridos de longitud $k$ no es, pues,
``parecido'' a una atención: \emph{es} una atención cuyo soporte el álgebra de
caminos determina exactamente (Proposición~\ref{prop:count}) y cuyos pesos
están congelados. Esto dicta el diseño: conservar el soporte algebraico y
reemplazar los pesos congelados de multiplicidad por pesos consulta--clave
aprendidos---el resultado es Walk Attention, la actualización aprendida
$\varphi$ del sustrato común de la Sección~\ref{sec:ground}.

% ---------------------------------------------------------------------
\subsection{La capa de Walk Attention}\label{sec:layer}
% ---------------------------------------------------------------------

Fijemos un alcance $k$ con soporte $\mathcal{S}_k$. Una capa de Walk Attention
con $H$ cabezas de dimensión $C$ calcula consultas, claves y valores con mapas
lineales, $Q^h=xW^h_Q$, $K^h=xW^h_K$, $V^h=xW^h_V$, y atiende \emph{solo sobre
los pares alcanzables} $(j,i)\in\mathcal{S}_k$:
\begin{equation}\label{eq:wa}
  \ell^h_{ji} = \tfrac{1}{\sqrt{C}}\langle Q^h_j,K^h_i\rangle,\qquad
  \alpha^h_{ji} = \softmax_{i:(j,i)\in\mathcal{S}_k}\ell^h_{ji},\qquad
  z^h_j = \textstyle\sum_{i:(j,i)\in\mathcal{S}_k}\alpha^h_{ji}V^h_i.
\end{equation}
El softmax se toma, para cada objetivo $j$, exactamente sobre las fuentes que
lo alcanzan por un recorrido de longitud $k$---un softmax \emph{segmentado}
sobre el soporte, nunca sobre los $n$ nodos. Las salidas de las cabezas se
concatenan y se añade un término ``raíz'' residual,
$\mathrm{WA}_k(x)_j=\Vert_h z^h_j + x_jW_{\!\text{root}}$. Como~\eqref{eq:wa}
solo toca pares de $\mathcal{S}_k$, la capa cuesta $O(|\mathcal{S}_k|HC)$ en
tiempo y memoria, no $O(n^2)$; los soportes se precalculan una vez por grafo
como los patrones de no-ceros de $\Adj^{\,k}$. La Figura~\ref{fig:layer}
muestra el flujo de datos y el Algoritmo~\ref{alg:wa} resume una capa.

\begin{figure}[t]
\centering
\begin{tikzpicture}[
  >={Stealth[round]}, node distance=6mm,
  box/.style={draw,rounded corners=2pt,minimum height=7mm,inner sep=3pt,font=\small,fill=blue!4},
  alg/.style={draw,rounded corners=2pt,minimum height=7mm,inner sep=3pt,font=\small,fill=green!6},
  lbl/.style={font=\scriptsize,gray}]
  \node[box] (x) {$x$};
  \node[box,right=8mm of x] (q) {$Q=xW_Q$};
  \node[box,below=2mm of q] (k) {$K=xW_K$};
  \node[box,below=2mm of k] (v) {$V=xW_V$};
  \node[alg,right=10mm of k] (att)
    {$\displaystyle \alpha_{ji}=\softmax_{i:(j,i)\in\mathcal{S}_k}\!\frac{\langle Q_j,K_i\rangle}{\sqrt C}$};
  \node[box,right=10mm of att] (agg) {$z_j=\sum_i \alpha_{ji}V_i$};
  \node[box,right=8mm of agg] (out) {$\mathrm{WA}_k(x)$};
  \node[alg,below=8mm of att,fill=green!10] (supp)
    {soporte $\mathcal{S}_k=\{(j,i):(\Adj^{\,k})_{ij}>0\}$};
  \draw[->] (x) -- (q.west);
  \draw[->] (x) |- (k.west);
  \draw[->] (x) |- (v.west);
  \draw[->] (q.east) -- (att.north west);
  \draw[->] (k.east) -- (att.west);
  \draw[->] (att.east) -- (agg.west);
  \draw[->] (v.east) -- ([yshift=-2mm]agg.west);
  \draw[->] (agg) -- (out);
  \draw[->,green!50!black] (supp.north) -- (att.south)
    node[midway,right,lbl,text=green!50!black] {restringe los pares};
  \draw[->,rounded corners] (x.north) -- ++(0,6mm) -| (out.north)
    node[pos=0.25,above,lbl] {raíz (residual)};
\end{tikzpicture}
\caption{Una capa de Walk Attention a alcance $k$. Consultas, claves y valores
son mapas lineales de los atributos; la atención se calcula \emph{solo} sobre
los pares alcanzables $\mathcal{S}_k$ que provee el álgebra de caminos (el
patrón de no-ceros de $\Adj^{\,k}$), se agrega y se suma un término raíz
residual. La caja verde del soporte distingue a Walk Attention de la atención
densa (todos los pares) y de la de un salto (aristas).}
\label{fig:layer}
\end{figure}

\begin{example}\label{ex:layer}
En el diamante del Ejemplo~\ref{ex:diamond} a alcance $k=2$ el soporte es el
único par $\mathcal{S}_2=\{(4,1)\}$: el nodo $4$ atiende directamente al nodo
$1$, el softmax sobre un singulete da $\alpha_{41}=1$, y
$\mathrm{WA}_2(x)_4=Vx_1+x_4W_{\!\text{root}}$---los dos recorridos paralelos
se vuelven una interacción \emph{directa}. Si se añade una segunda fuente $1'$
con flechas hacia $2$ y $3$, el soporte pasa a
$\mathcal{S}_2=\{(4,1),(4,1')\}$: el softmax pondera ahora las dos fuentes por
el ajuste de sus claves con la consulta $Qx_4$, de modo que el nodo $4$ lee
ambas fuentes a alcance $2$ directamente \emph{y elige entre ellas}---mientras
que una red de un salto debe empujar ambas señales por los vértices compartidos
$2,3$: la compresión misma que el over-squashing nombra.
\end{example}

Una red de Walk Attention de profundidad $L$ apila una capa por alcance
$k=1,\dots,L$, de modo que capas sucesivas leen recorridos progresivamente más
largos; cada capa va seguida de normalización de capa, una no linealidad
\textsc{elu} y dropout, y un MLP final mapea la incrustación a las salidas. La
normalización de capa importa en la práctica: sin ella la atención aprendida
tiene alta varianza entre inicializaciones; con ella, junto a un warmup corto y
recorte de gradiente, el entrenamiento converge de forma confiable
(Sección~\ref{sec:exp}).

\begin{algorithm}[t]
\caption{Capa de Walk Attention a alcance $k$ (dispersa).}\label{alg:wa}
\begin{algorithmic}[1]
\Require atributos $x\in\RR^{n\times d_{\text{in}}}$; soporte
  $\mathcal{S}_k=\{(j,i):(\Adj^{\,k})_{ij}>0\}$
\State $Q,K,V \gets xW_Q,\,xW_K,\,xW_V$ \Comment{por cabeza; formas $n\times HC$}
\For{cada $(j,i)\in\mathcal{S}_k$ y cabeza $h$}
  \State $\ell^h_{ji}\gets \langle Q^h_j,K^h_i\rangle/\sqrt{C}$
\EndFor
\State $\alpha^h_{ji}\gets$ $\softmax$ segmentado sobre $\{i:(j,i)\in\mathcal{S}_k\}$
\State $z^h_j\gets \sum_i \alpha^h_{ji}V^h_i$ \Comment{scatter-add hacia los objetivos}
\State \Return $\bigl\Vert_h z^h \;+\; xW_{\!\text{root}}$
\end{algorithmic}
\end{algorithm}

La Tabla~\ref{tab:relation} sitúa a Walk Attention entre los operadores
relacionados según los dos ejes de~\eqref{eq:attn-generic}: GAT aprende pesos
pero propaga por $L$ capas apiladas---el régimen del over-squashing; un graph
Transformer paga $O(n^2)$ y descarta la estructura; el operador de recorrido
tiene el soporte correcto pero no selecciona. Walk Attention ocupa la celda
restante.

\begin{table}[t]
\centering
\caption{Walk Attention frente a operadores de agregación relacionados, según
los dos ejes de~\eqref{eq:attn-generic}: soporte y pesos.}
\label{tab:relation}
\begin{tabular}{lccc}
\toprule
\textbf{Operador} & \textbf{Soporte} & \textbf{Pesos} & \textbf{Costo} \\
\midrule
GCN / paso de mensajes       & 1 salto      & fijos (grado) & $O(|Q_1|)$ \\
Atención sobre grafos (GAT)  & 1 salto      & aprendidos    & $O(|Q_1|)$ \\
Operador de recorrido $W^{(k)}$ & alcanzable por recorridos & fijos (multiplicidad) & $O(|\mathcal{S}_k|)$ \\
Transformer de grafo         & todos los pares & aprendidos & $O(n^2)$ \\
\textbf{Walk Attention (nuestra)} & \textbf{alcanzable por recorridos} & \textbf{aprendidos} & $O(|\mathcal{S}_k|)$ \\
\bottomrule
\end{tabular}
\end{table}

% ---------------------------------------------------------------------
\subsection{Un tratamiento formal sobre la familia de cuellos de botella}
\label{sec:theory}
% ---------------------------------------------------------------------

Sobre la familia de cadenas con cuello de botella de la Sección~\ref{sec:exp},
la separación entre el paso de mensajes de un salto y Walk Attention se puede
\emph{demostrar}, no solo medir. La construcción: $K$ fuentes llevan una
\emph{clave} one-hot---la asignación de claves es una permutación aleatoria
uniforme $\pi$---y un \emph{valor} one-hot; siguen $d$ capas de cuello de
botella $L_1,\dots,L_d$ de $M$ vértices intercambiables y sin atributos, con
capas consecutivas totalmente conectadas; un objetivo lleva una \emph{consulta}
one-hot $q$ y debe producir el valor de la fuente cuya clave es $q$ (etiqueta
$\pi^{-1}(q)$, azar $1/K$). Las redes tienen la profundidad estándar $L=d+1$
(una capa por salto, como en \cite{alon2021bottleneck}) y ancho oculto $w$.

\begin{proposition}\label{prop:collapse}
Para toda red de paso de mensajes cuya actualización
$h_v'=U\bigl(h_v,\mathrm{AGG}\{h_u:u\to v\}\bigr)$ se comparte entre vértices:
tras cada ronda, los $M$ vértices de cada capa de cuello de botella llevan
estados idénticos. Por tanto toda la influencia de las fuentes sobre el
objetivo pasa por un único estado $w$-dimensional por capa---mientras que
$KM^{d}$ recorridos la transportan (Proposición~\ref{prop:count}).
\end{proposition}

\begin{proof}
Inducción sobre la ronda: los vértices de una capa empiezan iguales (atributos
cero) y comparten su multiconjunto de entrada (todas las fuentes para $L_1$,
todo $L_{r-1}$ para $L_r$), así que $U$ y $\mathrm{AGG}$ compartidos los
mantienen iguales. \qed
\end{proof}

\begin{proposition}\label{prop:blind}
Si además los mensajes son lineales en el estado del emisor con coeficientes
independientes del contenido---como en GCN y GIN---entonces, con $L=d+1$ capas,
la salida del objetivo depende de las fuentes solo a través de $\sum_s x_s$,
que es la misma para toda $\pi$: las claves one-hot de una permutación suman el
vector de unos. La precisión esperada es exactamente $1/K$ para todo $d$, $w$ y
parámetros.
\end{proposition}

\begin{proof}
La información de las fuentes que entra a $L_1$ en la ronda $t$ necesita $d$
saltos más, llegando al objetivo en la ronda $t+d\le d+1$ solo si $t\le1$; en
la ronda $1$ el mensaje de la fuente $s$ es un mapa lineal fijo del atributo
crudo $x_s$ con un coeficiente compartido (las fuentes son isomorfas en el
grafo), así que el estado de $L_1$ es función de $\sum_s x_s$, cuyos bloques de
clave y de valor son independientes de $\pi$. Dada $q$, la etiqueta
$\pi^{-1}(q)$ es uniforme, así que toda predicción independiente de $\pi$
acierta $1/K$. \qed
\end{proof}

\begin{proposition}\label{prop:wa}
A alcance $k=d+1$ el soporte del objetivo es exactamente las $K$ fuentes (el
grafo es graduado; Proposición~\ref{prop:fns}). Con $W_Q$ proyectando la
casilla de consulta escalada por cualquier $\beta>0$, $W_K$ la casilla de
clave, $W_V$ la casilla de valor y $W_{\!\text{root}}=0$, la salida
$z_t=\sum_s\alpha_s e_{v(s)}$ de una sola capa de Walk Attention tiene su mayor
entrada en el valor correcto, en toda instancia, para todos $d$ y $M$.
\end{proposition}

\begin{proof}
El logit de la fuente $s$ es $\beta\,\mathbf{1}[\pi(s)=q]$, así que el softmax
pone su mayor peso $e^{\beta}/(e^{\beta}+K-1)$ en la única coincidencia
$s^\ast$; siendo los valores vectores distintos de la base, la mayor coordenada
de $z_t$ es $\alpha_{s^\ast}$, en $v(s^\ast)=\pi^{-1}(q)$. \qed
\end{proof}

\begin{example}\label{ex:blind}
Tómese $K=2$, $M=2$, $d=1$ y profundidad $L=2$. Las dos asignaciones de claves,
$\pi=\mathrm{id}$ y el intercambio, producen la \emph{misma} suma de fuentes
$\sum_s x_s$ (bloques de clave y de valor ambos $(1,1)$), así que por la
Proposición~\ref{prop:blind} una GCN o GIN calcula salidas idénticas en las dos
instancias---cuyas etiquetas difieren: precisión $1/2$, azar. Una capa de Walk
Attention a alcance $2$ atiende sobre $\{s_1,s_2\}$ con logits
$\beta\langle e_q,e_{\pi(s)}\rangle$, que distinguen las instancias desde la
primera multiplicación y seleccionan la fuente coincidente
(Proposición~\ref{prop:wa}).
\end{example}

\begin{remark}\label{rem:mechanisms}
Las proposiciones delimitan los mecanismos observados en la
Sección~\ref{sec:exp}. GAT escapa a la Proposición~\ref{prop:blind}---sus
coeficientes de la primera ronda dependen del contenido---pero debe hacer pasar
la asociación clave--valor por el canal de un solo vector de la
Proposición~\ref{prop:collapse}, re-codificado $d+1$ veces: queda por encima
del azar y se degrada con la profundidad. El operador de pesos fijos lee las
fuentes directamente pero con pesos \emph{iguales} (todas las multiplicidades
valen $M^d$), así que debe incrustar el multiconjunto completo en $w$
dimensiones en vez de seleccionar: se estanca entre ambos. Walk Attention lee
directamente \emph{y} selecciona (Proposición~\ref{prop:wa}), y resuelve la
tarea.
\end{remark}

% =====================================================================
\section{Experimentos}\label{sec:exp}
% =====================================================================

Evaluamos sobre una familia controlada con over-squashing severo y ajustable, y
luego sobre dos benchmarks externos. Todo el código, las configuraciones y los
registros se publican \cite{repo_walkattention}.

\paragraph{Recuperación en cadena con cuello de botella.}
Instanciamos la familia de la Sección~\ref{sec:theory} con $K{=}5$, $M{=}4$ y
profundidades $d\in\{2,3\}$: la multiplicidad hacia el objetivo es $KM^d$
(Proposición~\ref{prop:count}), mientras que el soporte más profundo consta
solo de los $K$ pares fuente--objetivo ($5/196\approx2.6\%$ de los pares a
$d{=}2$, $5/324\approx1.5\%$ a $d{=}3$). Todos los modelos usan ancho $16$,
cuatro cabezas cuando corresponde, $d{+}1$ capas/alcances y AdamW con warmup de
diez épocas, decaimiento coseno y recorte de gradiente; Walk Attention añade
normalización de capa. Reportamos precisión de validación en el objetivo (la
tarea sintética no tiene partición de prueba separada), media $\pm$ desviación
sobre cinco semillas. La Tabla~\ref{tab:main} cae exactamente donde las
Proposiciones~\ref{prop:collapse}--\ref{prop:wa} la sitúan
(Observación~\ref{rem:mechanisms}): GCN y GIN quedan clavados en el azar y allí
se quedan empíricamente; GAT se degrada abruptamente con la profundidad; el
operador de recorrido cruza el cuello en un paso pero no puede seleccionar la
fuente coincidente y se estanca; \textbf{Walk Attention resuelve la tarea
exactamente} en ambas profundidades y todas las semillas. Sin normalización de
capa ni warmup su precisión osciló entre $0.44$ y $1.00$ según la semilla; los
estabilizadores eliminan la varianza por completo---un artefacto de
optimización, no una propiedad del operador.

\begin{table}[t]
\centering
\caption{Precisión en el objetivo (media $\pm$ desv.\ sobre 5 semillas) en la
recuperación con cuello de botella, $K{=}5$, $M{=}4$; azar $1/K=0.20$; la
respuesta está a $d{+}1$ saltos. GCN/GIN quedan en el azar como exige la
Proposición~\ref{prop:blind} (el residuo es sesgo de mejor época).}
\label{tab:main}
\begin{tabular}{lcc}
\toprule
\textbf{Model} & \textbf{depth $d=2$} & \textbf{depth $d=3$} \\
\midrule
GCN (one-hop, linear aggregation)  & $0.21\pm0.01$ & $0.22\pm0.01$ \\
GIN (one-hop, linear aggregation)  & $0.21\pm0.01$ & $0.21\pm0.01$ \\
GAT (one-hop attention)            & $0.45\pm0.22$ & $0.29\pm0.17$ \\
Walk operator (fixed weights)      & $0.62\pm0.12$ & $0.50\pm0.12$ \\
\textbf{Walk Attention (ours)}     & $\mathbf{1.000\pm0.000}$ & $\mathbf{1.000\pm0.000}$ \\
\bottomrule
\end{tabular}
\end{table}

\paragraph{RingTransfer.}
Para confirmar el efecto en un benchmark independiente usamos
\textsc{RingTransfer}
\cite{digiovanni2023oversquashing}: un ciclo de $n$ nodos donde una fuente debe
transmitir su etiqueta one-hot (cinco clases, azar $0.2$) al objetivo
diametralmente opuesto, a $n/2$ saltos por los \emph{dos arcos} del anillo.
Cada nodo tiene dos vecinos, así que el número de recorridos de longitud $n/2$
que entran al objetivo---lo que una red de un salto debe comprimir tras $n/2$
capas---crece como $2^{n/2}$, mientras solo los dos arcos portan la señal.
Todos los modelos usan ancho $32$, la misma receta de optimización y $n/2$
capas/alcances. La Tabla~\ref{tab:ring} reporta los resultados: los anillos
pequeños son fáciles para todos; desde $n{=}14$ (distancia $7$) las líneas base
de un salto se degradan y se vuelven muy dependientes de la semilla---en
$n{=}18$ GAT cae a $0.29\pm0.19$ y GCN a $0.53\pm0.33$---mientras
\textbf{Walk Attention se mantiene en $1.000\pm0.000$}. La brecha se amplía con
el tamaño del anillo, esto es, con la severidad del over-squashing, como se
predijo.

\begin{table}[t]
\centering
\caption{Precisión en \textsc{RingTransfer} frente al tamaño $n$ (distancia
$n/2$), media $\pm$ desv.\ sobre 5 semillas; azar $0.20$.}
\label{tab:ring}
\begin{tabular}{lcccc}
\toprule
\textbf{Model} & $n{=}6$ & $n{=}10$ & $n{=}14$ & $n{=}18$ \\
\midrule
GCN                    & $1.00\pm0.00$ & $1.00\pm0.00$ & $0.81\pm0.13$ & $0.53\pm0.33$ \\
GAT                    & $1.00\pm0.00$ & $1.00\pm0.00$ & $0.69\pm0.29$ & $0.29\pm0.19$ \\
\textbf{Walk Attention (ours)} & $\mathbf{1.00\pm0.00}$ & $\mathbf{1.00\pm0.00}$ & $\mathbf{1.00\pm0.00}$ & $\mathbf{1.00\pm0.00}$ \\
\bottomrule
\end{tabular}
\end{table}

\paragraph{Colapsar recorridos perjudica.}
El cociente $\kQ/I$ de la Sección~\ref{sec:algebra} sugiere un remedio
alternativo: reemplazar las multiplicidades $n_k(i,j)$ por las dimensiones
menores del cociente, deduplicando recorridos equivalentes antes de agregar.
Comparamos ambos operadores en una ablación que cambia \emph{únicamente} el
operador, con datos, arquitectura y receta de la Tabla~\ref{tab:main} fijos. La
Tabla~\ref{tab:quotient} muestra al cociente claramente por debajo del operador
crudo en ambas profundidades: colapsar recorridos paralelos descarta la
multiplicidad de fuentes que la recuperación necesita. En este régimen los
caminos redundantes son \emph{señal}, no ruido; la contribución del álgebra es
el soporte, no la compresión. Si existe un régimen donde la compresión del
cociente ayude---caminos paralelos genuinamente redundantes, o clases de
caminos enrutadas a cabezas distintas---queda abierto.

\begin{table}[t]
\centering
\caption{Deduplicar recorridos vía $\kQ/I$ \emph{perjudica}: media $\pm$ desv.\
sobre 5 semillas, mismo protocolo de la Tabla~\ref{tab:main} (la fila del
operador coincide).}
\label{tab:quotient}
\begin{tabular}{lcc}
\toprule
\textbf{Operator} & \textbf{depth $d=2$} & \textbf{depth $d=3$} \\
\midrule
Quotient $\kQ/I$ (collapsed walks) & $0.39\pm0.04$ & $0.30\pm0.09$ \\
Walk operator (raw walks)          & $\mathbf{0.62\pm0.12}$ & $\mathbf{0.50\pm0.12}$ \\
\bottomrule
\end{tabular}
\end{table}

\paragraph{LRGB Peptides-func.}
Para evaluar con datos reales usamos Peptides-func \cite{dwivedi2022lrgb}
($10{,}873/2{,}331/2{,}331$ grafos moleculares, mediana $138$ nodos;
clasificación multietiqueta, Average Precision de prueba), con el mismo
embedding de átomos, cuatro capas, ancho $80$, pooling de media y AdamW para
todos los modelos; Walk Attention usa alcances $k=1,\dots,4$, cuyo soporte
sigue disperso ($1$--$8\%$ de $n^2$ a alcance $3$). Tabla~\ref{tab:lrgb}: Walk
Attention supera a GCN y es estadísticamente indistinguible de GAT
($0.526\pm0.007$ frente a $0.526\pm0.004$ sobre cuatro semillas). Dos
salvedades enmarcan la comparación. Primera, todos los modelos comparten un
presupuesto deliberadamente pequeño, así que los números absolutos quedan por
debajo de los publicados (GCN $0.593$ bajo el protocolo afinado de $500$k
parámetros \cite{dwivedi2022lrgb}, mayores aún en re-evaluaciones
\cite{tonshoff2023where}): la comparación es \emph{interna}, aislando el
operador de agregación. Segunda, el empate no se debe a la ausencia de caminos
paralelos (los péptidos contienen cientos de pares unidos por varios recorridos
cortos) sino a que el over-squashing es \emph{leve} en estos grafos pequeños.
La ventaja del soporte multi-salto es por tanto \emph{dependiente del régimen}:
decisiva cuando el squashing es severo
(Tablas~\ref{tab:main}--\ref{tab:ring}), modesta cuando es leve, a alrededor de
$1.6\times$ el tiempo de pared de GAT---muy por debajo de la atención densa.

\begin{table}[t]
\centering
\caption{LRGB Peptides-func, AP de prueba (media $\pm$ desv.\ sobre 4
semillas), presupuesto pequeño compartido que aísla el operador; los resultados
publicados bajo el protocolo estándar de $500$k son mayores para toda
arquitectura \cite{dwivedi2022lrgb,tonshoff2023where}.}
\label{tab:lrgb}
\begin{tabular}{lc}
\toprule
\textbf{Model} & \textbf{test AP} \\
\midrule
GCN                          & $0.492\pm0.011$ \\
GAT                          & $0.526\pm0.004$ \\
Walk Attention (ours)        & $0.526\pm0.007$ \\
\bottomrule
\end{tabular}
\end{table}

% =====================================================================
\section{Conclusión y Trabajo Futuro}\label{sec:conclusion}
% =====================================================================

Introdujimos Walk Attention, una arquitectura que mitiga el over-squashing
atendiendo sobre un soporte multi-salto provisto por el álgebra de caminos del
quiver subyacente. La Proposición~\ref{prop:count} da el soporte y sus
multiplicidades exactamente; las
Proposiciones~\ref{prop:collapse}--\ref{prop:wa} demuestran la separación en la
familia de cuellos de botella---el paso de mensajes con primera agregación
lineal queda en el azar para toda profundidad y ancho, mientras una capa de
Walk Attention en el alcance adecuado es exacta---y los experimentos la
confirman, con ventaja dependiente del régimen en \textsc{RingTransfer} y
Peptides-func. Un resultado negativo complementario mostró que comprimir
recorridos equivalentes vía el cociente es contraproducente cuando la
multiplicidad porta señal: el papel del álgebra es determinar el soporte de la
atención, no descartar información. La construcción es deliberadamente
elemental y autocontenida, una puerta de entrada de la teoría de
representaciones de quivers al aprendizaje de representaciones en grafos; es la
lectura aprendida de un sustrato común cuya lectura dinámica es el marco AIQ
\cite{zainea2026aiq}. Siguen varias direcciones.

\begin{itemize}
  \item \textbf{De los pesos aprendidos a la dinámica.} Inyectar los pesos de
  atención aprendidos en un AIQ como regla de evolución guiada por datos, y
  recíprocamente dejar que los análisis de estabilidad del AIQ interpreten lo
  aprendido---para nosotros el puente más prometedor.
  \item \textbf{Benchmarks con over-squashing severo.} Tareas reales de largo
  alcance con cuellos de botella agudos, comparadas contra rewiring
  \cite{topping2022understanding}, métodos multi-salto \cite{gutteridge2023drew}
  y graph Transformers.
  \item \textbf{Soportes con forma de cociente.} Usar $\kQ/I$ para
  \emph{particionar} los recorridos en clases enrutadas a cabezas de atención
  distintas, en lugar de comprimir; identificar qué relaciones $I$ sirven, y
  aprenderlas de los datos, es un problema atractivo.
  \item \textbf{Escalar el soporte.} Truncar, muestrear o aprender soportes
  admisibles más dispersos en grafos densos, donde el soporte de alcance $k$
  puede crecer.
\end{itemize}

\subsubsection*{Reproducibilidad.}
Los generadores de grafos, los operadores de recorrido, la capa Walk Attention,
las líneas base y los registros de experimentos se publican
\cite{repo_walkattention}; los cálculos del álgebra usan la biblioteca abierta
\texttt{aiq-quivers} \cite{repo_aiqquivers}.

% =====================================================================
\bibliographystyle{splncs04}
\bibliography{walk_attention}
% =====================================================================

\end{document}
